\documentclass[12pt]{article}
\usepackage[margin=2.2in]{geometry}
\usepackage{amsmath,amsthm,amssymb,mathtools}
\usepackage{setspace}
\setstretch{1.2}

\newtheorem*{theorem}{Theorem}
\newcommand{\N}{\mathbb{N}}

\begin{document}

\begin{theorem}[Infinitude of the primes]
  For all $n \in \N$,
  there exists $p \in \N$ such that
  $p > n$ and $p$ is prime.
\end{theorem}

\begin{proof}
  %
  Take $n \in \N$, write $m \coloneqq n! + 1$,
  and let $p$ be a prime factor of $m$.
  %
  \begin{itemize}
    %
    \item Note that $m \neq 1$ because $n! \neq 0$.
    \item Therefore, $p$ is prime.
    \item Suppose for contradiction that we do not have $p > n$.
    \item Then $p \leq n$.
    \item We have $p \mid m$ by definition of $p$.
    \item Note that $p > 0$ because $p$ is a factor of $m$.
    \item Also, $p \mid n!$ because $p \leq n$ by assumption and $p > 0$.
    \item So $p \mid 1$ by modular arithmetic.
    \item Thus, $p = 1$.
    \item But $1$ is not prime \textreferencemark.
      \qedhere
    %
  \end{itemize}
  %
\end{proof}

\end{document}
